% ============================================================
% 分片：s04  对应原文 PDF 第 61--66 页（第 66 页仅含附录 A 的收尾三段）
% 负责章节：§4 Concluding remarks；§5 Acknowledgments；\appendix；附录 A + A.1 + A.2
% 译者：translator-06   审校状态：初稿
% ============================================================

% ==== 原文 PDF 第 61 页 ====
\section{结语}

\begin{keytakeaway}
\textbf{核心物理图像：}扭变与上同调把物理系统中受约束的部分（真空流形、模空间及其形变）组织成可计算的代数结构；镜像对称等强有力几何等价性由此成为可精确陈述、可检验的命题。
\textbf{数学工具：}幂零微分诱导的上同调（BRST／BV、超对称扭变）、高阶代数与同伦代数（$A_\infty$、$L_\infty$、Poisson 顶点代数）、等变上同调与超对称局域化、形变理论和导出代数几何。
\textbf{主要结论：}弦对偶之网仍缺少在耦合 $g\sim1$ 处普遍有效的描述；更深地理解弦论与 M 理论的数学（尤其是高阶代数与导出几何）是走向强耦合统一图像的主要途径。
\end{keytakeaway}

本讲义对弦对偶中的数学所作的旋风式巡礼到此结束。遗憾的是，我们对这一主题的巡礼颇为简短，许多重要论题未能涉及。尽管如此，我们希望所选论题能为读者在这一充满活力、多面的领域中进一步积累专业知识提供有用的基础。

下面作简短的总结性评论，并回顾若干重要经验。当我们有一个自然的幂零微分时——正如在规范理论的（BV-）BRST 量子化中以及超对称理论的（扭变）中那样——上同调为分解与研究物理与数学系统的子空间或商提供了强有力的协变手段。在弦论语境中，世界面上的扭变产生拓扑弦，而拓扑弦一直是镜像对称等强有力几何等价性的来源。母链复形与上同调群本身都可以携带丰富而精细的代数结构，把这些受约束（或简化）的物理系统的算子（或态，本讲义未着重讨论）彼此联系起来。这些代数通常是高阶代数，同伦代数与 Poisson 顶点代数就是其中两个具体而普遍的实例。
% [Tutor Note: 物理图像：本讲义的方法论主线是“用幂零微分筛选受保护信息”——上同调丢弃 $Q$-恰当的自由度，只留下与真空流形、模空间几何对应的数据；而链层面的高阶运算则记录了被丢弃的形变信息如何组织成同伦代数。]

尽管我们只在几处简短评论中提及 BV-BRST 形式，仍建议读者参阅 [135--139] 中的详尽综述。在那里，上同调与同伦代数同样不可避免，即使未必出现任何超对称。它们还与导出代数几何（derived algebraic geometry）有着由来已久而优美的数学联系。

我们还简要提及了等变上同调；超对称局域化可视为这一概念的无穷维推广，可用于计算大量超对称可观测量，包括 BPS 态的关联函数。BPS 态可由扭变程序分离出来，它们还与模空间的几何深刻相关，支配物理理论的形变或不等价真空的族。确实，物理 Hilbert 空间的这些特殊子空间可以张成模空间的切空间。形变理论又以多种方式与高阶代数相联系，其中一些方式至今仍只是部分地被理解。

最后我们指出，弦对偶之网广袤而精细，贯穿紧化、微扰与非微扰动力学以及数学的许多分支。尽管许多强--弱对偶使我们得以触及弦论景观的各个角落，我们这些升级版的“傅里叶变换”在耦合 $g\sim 1$ 时便失去用处，因为不存在使耦合变弱的对偶框架。对于远离某个对偶框架中微扰极限的弦论，我们仍缺乏有用的描述；我们或许可以期望，对弦论与 M 理论数学的更深入理解能帮助我们抵达那里。

\section{致谢}

感谢 K. Costello、G. Moore 与 B. Williams 对本文初稿提出的有益意见，也感谢 M. Duff 与 W. Gu 对本文早期版本的意见。我们同样感谢 K. Costello、T. Dimofte、B. Williams 与 I. Saberi 的出色合作，这些合作塑造并澄清了我们对本讲义所涉论题的认识。NP 感谢 TASI 2021 的组织者提供作这些报告的机会并鼓励
% ==== 原文 PDF 第 62 页 ====
撰写本文，也感谢学生们的参与和富有洞见的问题。NG 与 NP 的工作由 University of Washington 支持；NP 还获得 DOE 资助 DE-SC0022347。

\appendix

\section{局域化概要}

\begin{keytakeaway}
\textbf{核心物理图像：}超对称路径积分可在不改变精确结果的前提下作 $Q$-恰当形变，使涨落被压制到 BPS 构型附近，从而把无穷维路径积分约化为低维（极端情形下有限维）积分乘以单圈行列式。
\textbf{数学工具：}等变上同调（$d_\sigma$、$\mathcal{L}_V$、Killing 矢量）、Duistermaat--Heckman 与 Atiyah--Bott 局域化、Stokes 定理、鞍点展开与单圈精确性、超行列式（SDet）。
\textbf{主要结论：}局域化公式把期望值写成局域化轨迹上的积分乘以法向涨落的单圈行列式的倒数；局域化项的选择等价于选取可观测量的对偶表示。
\end{keytakeaway}

超对称 QFT 往往具有可精确计算的量，且这些量在理论的形变下异常稳健。计算这些量的能力常常来自超对称局域化（supersymmetric localization，或简称局域化）这一不可或缺的工具。关于 QFT 中局域化的多种教学讲义，参见 [24, 140]。鉴于局域化在检验对偶与计算特殊类可观测量（包括扭变理论中）的关联函数方面的重要性，我们在此勾勒其基本思想，以备读者此前并不熟悉。

用局域化技术算出的众多量（各种（弯曲）空间上的配分函数、超对称指标、关联函数）往往可丰富地翻译为数学量。最明显的是，例如把同一配分函数的两种对偶表示算出并令二者相等，就能得到高度非平凡的积分恒等式。更深刻地说，超对称局域化给出了等变上同调中经典局域化概念的物理实现／推广，后者归功于 Duistermaat--Heckman [141] 与 Atiyah--Bott [142]。

我们将先用几何语言引入这一思想，再把基本思想翻译成超对称场论的语言。如生活中一贯的情形，我们寻求计算路径积分的方法，最好能精确计算。若理论是自由的，积分是高斯型的，当然可以精确算出。在弱耦合理论中，我们可以希望把路径积分按微扰级数展开；但即使能够对自由理论附近的微扰级数求和，这至多是一个在耦合常数上收敛半径为零的渐近级数。\footnote{不过应注意，resurgence 这一主题旨在从路径积分的微扰级数推知其非微扰部分，已被证明在许多物理与数学应用中极为强大。它与 Stokes 现象、墙穿越（wall-crossing）等密切相关。入门可参见 [143, 144]。}此外，在研究对偶时，弱耦合理论的对偶框架几乎必然是强耦合的，不适于这类分析。然而在存在超对称时，路径积分中的玻色子与费米子贡献彼此抵消，这种抵消可用来把计算约化为更低维场论的路径积分，甚至更好，约化为在由适当超对称变换固定的有限维场空间上的积分：我们说路径积分局域化到 BPS 场构型。

\subsection{等变上同调中的局域化}

作为热身，先回顾等变上同调中局域化的思想：利用积分流形上的对称性精确计算积分。想法是计算定义在带对称群 $G$ 的流形 $\mathcal{M}$ 上的积分。（请留意这与存在规范对称性时的路径积分问题的形式相似性！）可以改为尝试计算 $\mathcal{M}/G$ 上的积分，即每条轨道只取一个代表。当 $G$ 在 $\mathcal{M}$ 上自由作用时这很好用，但当 $G$ 对 $\mathcal{M}$ 的作用有不动点轨迹时，
% ==== 原文 PDF 第 63 页 ====
该空间可能奇异。与其设法为（可能奇异的）空间 $\mathcal{M}/G$ 上的积分赋予意义，不如考虑 $\mathcal{M}$ 上 $G$-等变形式的积分，即 $G$-等变上同调的元素。我们关于局域化的简略讨论大体遵循优美的导论 [145]，建议读者参阅该文以获取更细致的处理与细节。

为简单起见，考虑闭、连通、有向 Riemann 流形 $\mathcal{M}$ 上的 $U(1)$ 作用。设生成该 $U(1)$ 作用的矢量场 $V$ 是 Killing 矢量场：$\mathcal{L}_V g_{\mu\nu}=0$。$U(1)$-等变上同调类 $[\omega]$ 由满足 $\mathcal{L}_V\omega=0$ 的 $U(1)$-不变微分形式（通常具有非齐次形式次数，并被等变参数 $\sigma$ 的多项式修饰）代表，且该形式等变闭 $d_\sigma\omega=0$，模去等变恰当形式的添加 $\omega\sim\omega+d_\sigma\lambda$（其中 $\mathcal{L}_V\lambda=0$）。当对称性自由作用时，这些上同调群与通常上同调一致；否则它们为非自由作用给出了良定义的上同调概念。

$U(1)$-等变形式的积分照常来自对最高次形式部分的积分，即形式次数等于 $\mathcal{M}$ 的维数的部分。若要积分的等变恰当形式为 $\omega=d_\sigma\lambda$，则其最高次形式分量自动为 de Rham 恰当的 $\omega|_{\Omega^n(\mathcal{M})}=d\lambda|_{\Omega^{n-1}(\mathcal{M})}$，故其积分由 Stokes 定理为零：
% [存疑: 原文以 $\omega|_{\Omega^n(\mathcal{M})}$ 表示“取 $\omega$ 的 $n$ 形式分量”，该记号并不常见；此处照抄原样以便核对。]
\begin{equation}\tag{A.1}
\int_{\mathcal{M}}\omega = \int_{\mathcal{M}} d\lambda|_{\Omega^{n-1}(\mathcal{M})} = 0,
\end{equation}
因此 $\mathcal{M}$ 上等变闭形式的积分只依赖于其等变上同调类。

等变积分的优美之处在于它们会局域化。这意味着，若在 $\mathcal{M}$ 上积分一个 $U(1)$-等变闭形式，积分只会从 $\mathcal{M}$ 上的 $U(1)$ 不动点获得贡献。下面快速勾勒该命题的证明，它将容易推广到超对称场论的情形。
% [Tutor Note: 推导补足：这里的关键与 (A.4) 的 $t\to\infty$ 极限相同——等变闭性保证积分与 $t$ 无关，而 $t\to\infty$ 时因子 $e^{\sigma t|V|^2}$ 把测度压成集中于 $V=0$ 的 delta 函数，故只有不动点（Killing 矢量的零点）贡献。]

由于积分只依赖于形式的 $U(1)$-等变上同调类，我们可以自由地把 $\omega$ 换成更方便的代表。具体而言，用 $d_\sigma$-恰当项形变被积式：
\begin{equation}\tag{A.2}
\omega \mapsto \omega_t := \omega\, e^{t\,d_\sigma\lambda},\qquad t\in\mathbb{R},\ \ \mathcal{L}_V\lambda=0.
\end{equation}
按构造，改变 $t$ 不改变被积式的上同调类，因为 $\omega$ 与 $\omega_t$ 等变上同调：
\begin{equation}\tag{A.3}
\omega_t = \omega\left(1+t\,d_\sigma\lambda+\ldots\right) = \omega + d_\sigma\left((-1)^{F(\omega)}t\,\omega\lambda+\ldots\right)
\end{equation}
于是，由于 $\mathcal{M}$ 没有边界，原积分对 $t$ 的导数也为零。这意味着可以在任何方便的 $t$ 值处计算积分，所得结果等价于 $t=0$ 时的原积分。可以把 $t=\frac{1}{\hbar}$ 看作某个辅助量子参数的倒数，于是 $t\to\infty$ 是半经典极限。$t$ 无关性告诉我们：只需在鞍点附近作单圈近似展开，就能得到原积分的精确结果！

若 $d_\sigma\lambda$ 的 0 形式部分负半定且最大值等于 0，则该极限存在。设情形如此，则在 $t\to\infty$ 极限下积分由 $-d_\sigma\lambda$ 的零点（即
% ==== 原文 PDF 第 64 页 ====
极小点）支配。进一步取 $\lambda$ 为与矢量场 $V$ 对偶的 1-形式 $\mu$，于是 $d_\sigma\mu = d\mu+\sigma|V|^2$。此时我们在计算积分
\begin{equation}\tag{A.4}
\lim_{t\to\infty}\int_{\mathcal{M}}\omega\, e^{t\,d\mu} e^{\sigma t|V|^2}.
\end{equation}
由于 $e^{t\,d\mu}$ 是 2-形式的指数，可展开为形式的多项式，其次数不高于 $n/2$，其中 $n$ 是流形的维数（更高次形式的积分都会平凡地为零）。另一因子 $e^{\sigma t|V|^2}$ 对 $t$ 有指数依赖，会压过第一项，并在 $t\to\infty$ 极限下趋于一个 delta 函数，其支集为 Killing 矢量 $V$ 的零点轨迹。

取上述 $\lambda$ 后我们看到：半经典极限在 Killing 矢量的零点轨迹上有 delta 函数支撑，至多相差一个适当的 Jacobian 因子，它度量 $V$ 的零点轨迹附近的一阶涨落。若流形上有多个 Killing 矢量使被积式等变闭，我们可以选择对哪个矢量或其组合作局域化。当然，我们一般也有选取 $\lambda$ 的自由，只需满足“半经典”极限存在这一条件。

\subsection{量子场论中的局域化}

下面把上述论证推广到量子场论。设存在费米子对称性生成元 $Q$，其平方为玻色子对称性生成元 $Q^2=J$；$Q$ 类比于 $U(1)$-等变外微分 $d_\sigma$，$J$ 类比于生成该对称性的 Lie 导数 $\mathcal{L}_V$。在最简单的情形中 $J=0$，$Q$ 就是幂零的费米子对称性。更一般地，$J$ 可以是时空对称性、整体对称性或规范对称性的某种线性组合。若 $Q$ 是 BRST 荷，则下面的局域化论证给出路径积分的规范固定；若它是通常的超荷，则得到超对称局域化的一般处方。实践中，我们常关心的是 BRST 微分与超对称生成元的组合微分 $Q$，正如 2.1 节开头所讨论的。
% [存疑: 原文 "as discussed in at the beginning of Section 2.1" 中 "in at" 多一介词，疑为笔误，译文按 "as discussed at the beginning of Section 2.1" 理解。]
% [Tutor Note: 易错预警：$Q^2=J$ 与 $d_\sigma$、$\mathcal{L}_V$ 的对应是 $Q\leftrightarrow d_\sigma$、$J\leftrightarrow\mathcal{L}_V$ 的形式类比，并不是超对称代数的定理；在纯 BRST 情形 $J=0$（$Q$ 幂零），而在超对称局域化中 $J$ 通常是 $R$ 对称性、时空旋转等玻色对称性的组合。]

与上述类比，我们预期超对称量子场论的路径积分会局域化到被玻色对称性 $J$ 固定的场构型上。当然，我们是在非紧的无穷维 $\mathcal{M}$ 上积分，实践中需要考虑各种微妙之处。这里只勾勒主要思想。

我们的目标是计算如下形式的积分
\begin{equation}\tag{A.5}
\langle O\rangle = \int_{\mathcal{F}}\mathcal{D}\phi\, O\, e^{-S[\phi]},
\end{equation}
其中 $\mathcal{F}$ 是场空间，$O$ 是 $Q$-闭、$J$-不变的可观测量。它可以是局域算子、一组局域算子，或具有非平凡空间延展的算子，如 Wilson 线。我们始终假设 $S[\phi]$ 是 $Q$-闭且 $J$-不变的。我们还假设 $Q$、$J$ 无反常，从而测度在 $Q$ 作用下也不变。由此，若 $O$ 是 $Q$-恰当的，即 $O=[Q,O']$，则该期望值为零；这是
% ==== 原文 PDF 第 65 页 ====
Stokes 定理的无穷维推广：
% [Tutor Note: 推导补足：这一步只是把 $[Q,O']$ 写成对称变分 $\delta_Q O'$ 并移到整个被积式上，得到场空间中的全导数。其成立依赖两个前提：$Q$ 不变测度（无反常）与场空间无边界（被积式在无穷远处衰减足够快），后者见脚注 54。]
\begin{align}
\langle[Q,O']\rangle &= \int_{\mathcal{F}}\mathcal{D}\phi\,[Q,O']e^{-S[\phi]} \tag{A.6}\\
&= \int_{\mathcal{F}}\mathcal{D}\phi\,\delta_Q\left(O'e^{-S[\phi]}\right) = 0. \tag{A.7}
\end{align}
我们用到了这样的事实：$Q$ 的作用是对称变分 $\delta_Q$，可以移到整个被积式上，成为场空间中的全导数。\footnote{这一论证默认了由此产生的边界项为零，即被积式在场空间趋于无穷时衰减足够快。后文均作此假设；不过也存在例外，会引出有趣的微妙之处。}

我们再次处于这样的情形：要计算的积分只依赖于被积对象的上同调类，此处即 $O$ 的 $Q$-上同调类。下面重复与前面同类的局域化论证。与简单的几何例子不同，由于场空间 $\mathcal{F}$ 的无穷维性，积分的分析要微妙得多。我们假设路径积分没有红外发散；实践中，这往往需要开启 $\Omega$-背景，或把超对称理论置于紧流形上。

与前面一样，用 $Q$-恰当项形变积分，使积分最终局域化到场空间中的固定轨迹：
\begin{equation}\tag{A.8}
\langle O\rangle_t = \int_{\mathcal{F}}\mathcal{D}\phi\, O\, e^{-S[\phi]-t[Q,V[\phi]]}.
\end{equation}
$V[\phi]$ 必然是费米子的，且必须 $J$-不变，才能构成良好的等变上同调问题。我们取 $[Q,V[\phi]]$ 的玻色部分为半正定，使得 $t>0$ 时指数有下界。在各项假设成立的前提下，$t$ 无关性的论证与前一情形完全相同。我们将利用这一 $t$ 无关性取 $t\to\infty$ 极限，此时积分由形变项的鞍点支配。同样，我们有选取 $V[\phi]$ 的自由，而这种自由可以揭示积分的不同对偶表示，即在不同外观的自由度上作局域化。在 BRST 语境中，我们有按需作规范固定的自由，不同的规范往往带来不同的洞见。当然，物理可观测量最终与规范选取无关，BPS 可观测量同样必须在不同的局域化项下给出相同答案。

为计算该量，把全部场在其鞍点值附近展开，$\phi=\phi_0+\frac{1}{\sqrt{t}}\delta\phi$，于是作用量成为 $S[\phi_0]+\frac{1}{2}\iint\frac{\delta^2[Q,V[\phi]]}{\delta\phi^2}\big|_{\phi=\phi_0}(\delta\phi)^2$。在严格的 $t\to\infty$ 极限下，该展开在参数 $\frac{1}{t}$ 上单圈精确。剩下的只是对 $\delta\phi$ 中垂直于局域化轨迹 $\mathcal{F}_Q$ 的涨落积分，得到玻色子与费米子场的单圈行列式。结果因此为
% [Tutor Note: 推导补足：形变项按 $\frac{1}{\sqrt t}\delta\phi$ 展开后，二次项系数即 $\frac{\delta^2[Q,V]}{\delta\phi^2}$，故涨落积分是高斯型（费米子部分为 Berezin 积分），给出“超行列式”因子；由于 $Q$ 恰当形变不改变积分，这个只保留到二次阶的结果竟然是精确的——这正是“单圈精确”的含义。]
\begin{equation}\tag{A.9}
\langle O\rangle = \int_{\mathcal{F}_Q}\mathcal{D}\phi_0\, O|_{\phi=\phi_0}\, e^{-S[\phi_0]}\left(\operatorname{SDet}\frac{\delta^2[Q,V[\phi_0]]}{\delta\phi_0^2}\right)^{-1}.
\end{equation}
% [背景补充: SDet 为超行列式，等于玻色行列式与费米子行列式之比；法向玻色与费米涨落的贡献正是以这种方式相互抵消，这是局域化结果如此简洁的根本原因。]

路径积分已等变地局域化为场空间 BPS 子轨迹 $\mathcal{F}_Q$ 上的一个低维积分。原始的经典作用量在 BPS 场构型上的取值，
% ==== 原文 PDF 第 66 页 ====
由概括该轨迹法向场涨落的单圈行列式修正。在具体语境中，这种局域化甚至可能把计算约化为普通的积分。选择局域化方案的自由，是获得各种可观测量不同对偶表示的有力手段。此外，选择费米子对称性 $Q$ 的自由，使得可以计算原理论中大量超对称量。在（拓扑）扭变场论的语境中，可以在局域化程序中使用扭变超荷来计算超对称指标，它们往往有有趣的数学表达式。

在本文强调的二维 $\mathcal{N}=(2,2)$ 理论语境中，可以容易地通过局域化计算它们的 Witten 指标与椭圆亏格等量 [146, 147]。这些理论在二维球面上的配分函数近来也已由局域化算出 [148]，并被用于直接从 $A$ 模型提取 Gromov--Witten 不变量而无需借助镜像对称 [149]。此外，如正文所勾勒的，拓扑环的关联函数可以用这些方法计算 [9, 20]。由于这些结果在别处已有充分讨论，且不参与本讲义余下的内容，我们仅满足于电报式的说明，并请读者参阅所引文献。

% [跨页接续: 原文第 65 页/第 66 页——附录 A.2 的末尾两段自 65 页续至 66 页上半，本分片按语义收全]
% [分片边界: 附录 A 的正文延伸到原文第 66 页上半；本分片止于其最后一段（"…referring the reader to the cited texts."），
%  s05 自同页的 "B. A-type Quantum Mechanics and Morse Theory" 标题起接续，二者装配时直接拼接、不重复。]
